% BUILD11. Descritivo.
\begin{table}[htbp]\centering
\caption{Income gradient in referral under competing-risks incidence and restricted mean time, by birth cohort and postcode income quartile}
\label{tab:E2}
\footnotesize
\setlength{\tabcolsep}{4pt}
\begin{tabular}{lrrrrrr}
\toprule
Cohort & ref Q1 & ref Q2 & ref Q3 & ref Q4 & RMST Q1 & RMST Q4 \\
\midrule
2010 & 0.530 & 0.520 & 0.566 & 0.652 & 222 & 199 \\
2011 & 0.505 & 0.483 & 0.530 & 0.613 & 252 & 234 \\
2012 & 0.607 & 0.596 & 0.625 & 0.701 & 240 & 222 \\
2013 & 0.722 & 0.722 & 0.760 & 0.815 & 212 & 182 \\
2014 & 0.734 & 0.729 & 0.757 & 0.828 & 222 & 199 \\
2015 & 0.790 & 0.782 & 0.812 & 0.876 & 199 & 179 \\
2016 & 0.855 & 0.853 & 0.877 & 0.923 & 184 & 167 \\
\bottomrule
\end{tabular}
\\[3pt]{\footnotesize \textit{Note:} Cumulative incidence of referral at 365 days (Aalen--Johansen, competing risk = aging out) by CEP income quartile (Q1 lowest), and restricted mean days in the eligible queue within 365 days for Q1 and Q4. Replaces the Kaplan--Meier wait tables (S4a, S4b), which censored aging out. A higher incidence and a lower RMST both mean faster referral.}
\end{table}
